\documentclass[12pt,a4paper,oneside]{book}

% =====================================================================
% PACOTES ESSENCIAIS
% =====================================================================
\usepackage[brazilian]{babel}
\usepackage[T1]{fontenc}
\usepackage[utf8]{inputenc}
\usepackage{amsmath}
\usepackage{amssymb}
\usepackage{amsthm}
\usepackage{geometry}
\usepackage{fancyhdr}
\usepackage{hyperref}
\usepackage{booktabs}
\usepackage{mathtools}

% =====================================================================
% GEOMETRIA
% =====================================================================
\geometry{
    left=2.5cm,
    right=2.5cm,
    top=3cm,
    bottom=3cm,
    headheight=14pt
}

% =====================================================================
% CABECALHO E RODAPE
% =====================================================================
\pagestyle{fancy}
\fancyhf{}
\fancyhead[L]{\textit{Capitulo 5: Comportamento Exponencial}}
\fancyhead[R]{\thepage}
\fancyfoot[C]{\small Principios de Modelagem Matematica}
\renewcommand{\headrulewidth}{0.4pt}
\renewcommand{\footrulewidth}{0.4pt}

% =====================================================================
% AMBIENTES
% =====================================================================
\theoremstyle{definition}
\newtheorem{definition}{Definicao}[chapter]
\newtheorem{example}[definition]{Exemplo}
\newtheorem{remark}[definition]{Observacao}
\newtheorem{problem}{Problema}

% =====================================================================
% DOCUMENTO
% =====================================================================
\title{\textbf{Capitulo 5: Comportamento Exponencial}\\[0.4cm]
\large Principios de Modelagem Matematica}
\author{Baseado em Dym, C. L.}
\date{2026-04-14}

\begin{document}

\maketitle
\tableofcontents
\newpage

\chapter{Comportamento Exponencial}

\section{Ideia central do capitulo}

Muitos processos fisicos, biologicos e economicos variam de forma proporcional ao proprio estado atual.
Quando isso acontece, aparece naturalmente o comportamento exponencial.

A ideia fundamental e:
\begin{equation}
\frac{dy}{dt} = ky,
\end{equation}
onde $k$ e uma constante.

A solucao e:
\begin{equation}
y(t)=y_0e^{kt}.
\end{equation}

\begin{remark}
Se $k>0$, temos crescimento exponencial. Se $k<0$, temos decaimento exponencial.
\end{remark}

\section{Modelo basico de crescimento e decaimento}

\subsection{Derivacao rapida}

Partindo de $dy/dt=ky$ com condicao inicial $y(0)=y_0$:
\begin{equation}
\frac{1}{y}dy = k\,dt
\end{equation}
\begin{equation}
\int \frac{1}{y}dy = \int k\,dt
\end{equation}
\begin{equation}
\ln|y| = kt + C
\end{equation}
\begin{equation}
y(t)=Ce^{kt}
\end{equation}

Usando $y(0)=y_0$, segue $C=y_0$.

\subsection{Constante de tempo e meia-vida}

Para decaimento ($k<0$), definimos $k=-\lambda$ com $\lambda>0$:
\begin{equation}
y(t)=y_0e^{-\lambda t}.
\end{equation}

Constante de tempo:
\begin{equation}
\tau=\frac{1}{\lambda}.
\end{equation}

Meia-vida $t_{1/2}$:
\begin{equation}
\frac{y_0}{2}=y_0e^{-\lambda t_{1/2}} \Rightarrow t_{1/2}=\frac{\ln 2}{\lambda}.
\end{equation}

\section{Exemplos classicos de modelagem}

\subsection{Exemplo 1: juros compostos continuos}

Se um capital $C(t)$ rende proporcionalmente ao proprio valor:
\begin{equation}
\frac{dC}{dt}=rC,
\end{equation}
entao
\begin{equation}
C(t)=C_0e^{rt}.
\end{equation}

Se $C_0=1000$ e $r=0.08$ ao ano, em 3 anos:
\begin{equation}
C(3)=1000e^{0.24}\approx 1271.25.
\end{equation}

\subsection{Exemplo 2: decaimento radioativo}

Numero de nucleos $N(t)$:
\begin{equation}
\frac{dN}{dt}=-\lambda N \Rightarrow N(t)=N_0e^{-\lambda t}.
\end{equation}

Se meia-vida e 5 dias, entao:
\begin{equation}
\lambda=\frac{\ln 2}{5}.
\end{equation}

Apos 15 dias:
\begin{equation}
N(15)=N_0e^{-\lambda 15}=N_0\left(\frac{1}{2}\right)^3=\frac{N_0}{8}.
\end{equation}

\subsection{Exemplo 3: resfriamento de Newton}

Modelo:
\begin{equation}
\frac{dT}{dt}=-k(T-T_a),
\end{equation}
onde $T_a$ e temperatura ambiente.

Definindo $\theta=T-T_a$:
\begin{equation}
\frac{d\theta}{dt}=-k\theta \Rightarrow \theta(t)=\theta_0e^{-kt}.
\end{equation}

Logo:
\begin{equation}
T(t)=T_a + (T_0-T_a)e^{-kt}.
\end{equation}

\subsection{Exemplo 4: carga e descarga de capacitor RC}

Para descarga de capacitor:
\begin{equation}
\frac{dV}{dt}=-\frac{1}{RC}V,
\end{equation}
com solucao:
\begin{equation}
V(t)=V_0e^{-t/(RC)}.
\end{equation}

Constante de tempo do circuito:
\begin{equation}
\tau=RC.
\end{equation}

\section{Interpretacao em escala logaritmica}

Tomando logaritmo de $y(t)=y_0e^{kt}$:
\begin{equation}
\ln y = \ln y_0 + kt.
\end{equation}

Isso e uma reta em grafico semilog ($\ln y$ vs $t$), com inclinacao $k$.

\begin{remark}
Em dados experimentais, linearizar em log ajuda a identificar rapidamente se o processo e exponencial.
\end{remark}

\section{Conexao com modelos nao lineares}

Mesmo quando um sistema e nao linear, o comportamento perto de equilibrio pode ser exponencial apos linearizacao.

Se $\dot{x}=f(x)$ e $x^*$ e equilibrio, para perturbacao $\eta=x-x^*$:
\begin{equation}
\dot{\eta}\approx f'(x^*)\eta.
\end{equation}

Assim:
\begin{equation}
\eta(t)\approx \eta_0e^{f'(x^*)t}.
\end{equation}

Sinal de $f'(x^*)$ decide estabilidade local.

\section{Exemplo completo de ajuste de parametro}

Suponha um processo de decaimento com dados:
\begin{center}
$t$ (h): $0, 2, 4$ \quad e \quad $y$: $80, 48.5, 29.4$
\end{center}

Modelo: $y(t)=y_0e^{-\lambda t}$ com $y_0=80$.

Usando $t=2$:
\begin{equation}
48.5 = 80e^{-2\lambda}
\end{equation}
\begin{equation}
e^{-2\lambda}=0.60625
\end{equation}
\begin{equation}
\lambda\approx -\frac{1}{2}\ln(0.60625)\approx 0.25\,\text{h}^{-1}.
\end{equation}

Previsao em $t=4$:
\begin{equation}
y(4)=80e^{-1}\approx 29.4,
\end{equation}
coerente com os dados.

\section{Modelo logistico de crescimento}

O modelo exponencial $dP/dt = rP$ preve crescimento ilimitado, o que nao e realista.
O modelo logistico, proposto por Verhulst (1838), inclui uma capacidade de suporte $K$:
\begin{equation}
\frac{dP}{dt} = rP\!\left(1 - \frac{P}{K}\right).
\end{equation}

Interpretacao dos termos:
\begin{itemize}
    \item Para $P \ll K$: $(1 - P/K) \approx 1$ e o crescimento e praticamente exponencial.
    \item Para $P \to K$: a taxa de crescimento tende a zero.
    \item Equilibrios: $P^* = 0$ (instavel) e $P^* = K$ (estavel).
\end{itemize}

\subsection{Solucao analitica pelo metodo de fracoes parciais}

Separando variaveis e usando fracoes parciais em $P(K-P)$:
\begin{equation}
\int \frac{dP}{P(K-P)} = \frac{1}{K}\int\left(\frac{1}{P}+\frac{1}{K-P}\right)dP = rt + C.
\end{equation}

Apos algebra e aplicando $P(0) = P_0$:
\begin{equation}
P(t) = \frac{K}{1 + \left(\dfrac{K - P_0}{P_0}\right)e^{-rt}}.
\end{equation}

Verificacao: $P(0) = P_0$ e $\lim_{t \to \infty} P(t) = K$. $\checkmark$

O ponto de inflexao (maxima taxa de crescimento) ocorre em $P = K/2$, no instante:
\begin{equation}
t^* = \frac{1}{r}\ln\!\left(\frac{K - P_0}{P_0}\right).
\end{equation}

\begin{remark}
O grafico de $P(t)$ tem forma de S (sigmoide): crescimento acelerado ate $K/2$, depois
desacelerado ate saturar em $K$. Esse padrao e chamado de curva logistica.
\end{remark}

\subsection{Exemplo numerico}

Seja $r = 0.3\,\text{dia}^{-1}$, $K = 1000$ e $P_0 = 50$.

Ponto de inflexao: $t^* = \frac{1}{0.3}\ln(19) \approx 9.8\,\text{dias}$.

Em $t = 10\,\text{dias}$:
\begin{equation}
P(10) = \frac{1000}{1 + 19\,e^{-3}} \approx \frac{1000}{1.949} \approx 513.
\end{equation}

\section{Equacoes de Lanchester}

Em conflitos militares, Lanchester (1916) propôs modelos para a evolucao de duas forcas
$E(t)$ e $F(t)$ em combate (aparece no Capitulo 5 do livro de Dym como aplicacao de EDOs
acopladas):
\begin{equation}
\frac{dE}{dt} = -\lambda_F\,F, \qquad \frac{dF}{dt} = -\lambda_E\,E,
\end{equation}
onde $\lambda_E$ e $\lambda_F$ sao as eficacias de combate de cada forca.

\subsection{Invariante de Lanchester}

Multiplicando a primeira equacao por $2\lambda_E E$ e a segunda por $2\lambda_F F$ e somando:
\begin{equation}
\frac{d}{dt}\!\left(\lambda_E E^2 - \lambda_F F^2\right) = 0.
\end{equation}

Portanto:
\begin{equation}
\lambda_E E^2 - \lambda_F F^2 = \lambda_E E_0^2 - \lambda_F F_0^2 = \text{constante}.
\end{equation}

Essa e a \textbf{Lei Quadratica de Lanchester}: permite determinar o vencedor sem resolver
o sistema completo. Se $\lambda_E E_0^2 > \lambda_F F_0^2$, a forca $E$ vence.

\begin{remark}
As equacoes de Lanchester mostram que tamanho quadra (lei quadratica), nao cresce linearmente.
Dobrar o tamanho de uma forca quadriplica seu efeito de combate, dada eficacia constante.
\end{remark}

\section{Tabela comparativa de modelos de crescimento e decaimento}

\begin{center}
\begin{tabular}{llll}
\toprule
\textbf{Modelo} & \textbf{EDO} & \textbf{Solucao} & \textbf{Comportamento} \\
\midrule
Crescimento puro  & $\dot{y}=ky$ & $y_0e^{kt}$ & crescimento ilimitado \\
Decaimento puro   & $\dot{y}=-\lambda y$ & $y_0e^{-\lambda t}$ & decai a zero \\
Newton resfriamento & $\dot{T}=-k(T-T_a)$ & $T_a+(T_0-T_a)e^{-kt}$ & satura em $T_a$ \\
Logistico (Verhulst) & $\dot{P}=rP(1-P/K)$ & sigmoide (secao~anterior) & satura em $K$ \\
Lanchester & $\dot{E}=-\lambda_F F$ & hiperbolico & um lado zera \\
RC (capacitor) & $\dot{V}=-V/(RC)$ & $V_0e^{-t/(RC)}$ & decai a zero \\
\bottomrule
\end{tabular}
\end{center}

\section{Roteiro pratico para identificar comportamento exponencial}

\begin{enumerate}
    \item Verifique se a taxa de variacao parece proporcional ao estado atual.
    \item Escreva a EDO na forma $dy/dt=ky$ (ou equivalente apos mudanca de variavel).
    \item Resolva e use condicao inicial para obter a constante.
    \item Use grafico semilog para validar linearidade.
    \item Estime parametros ($k$, $\lambda$, $\tau$, meia-vida) a partir dos dados.
\end{enumerate}

\section{Exercicios curtos}

\begin{problem}
Uma populacao cresce com taxa proporcional ao tamanho, com $P(0)=120$ e $k=0.03$ dia$^{-1}$. Encontre $P(20)$.
\end{problem}

\begin{problem}
Um material tem meia-vida de 8 horas. Qual fracao resta apos 24 horas?
\end{problem}

\begin{problem}
No modelo de resfriamento de Newton, derive a expressao para o tempo necessario para atingir uma temperatura alvo $T^*$.
\end{problem}

% =====================================================================
% EXEMPLO TIPO PROVA
% =====================================================================
\section{Exemplo tipo prova: datacao por carbono-14}

\begin{problem}[Prova -- IIo semestre 2013]
Suponha que uma amostra de fossil tenha sido datada por carbono-14 e o resultado
obtido tenha sido de $570.000$ anos. Dado: concentracao natural em organismos
vivos $= 1$ atomo em $10^{12}$; meia-vida do ${}^{14}$C $= 5700$ anos.
\begin{enumerate}
    \item Encontre a concentracao de ${}^{14}$C supostamente encontrada no osso.
    \item Se a concentracao e determinada com precisao de 1 parte em $10^{23}$,
    encontre a incerteza na estimativa da idade.
    \item Em funcao dos valores encontrados, critique a estimativa feita para
    a idade do osso.
\end{enumerate}
\end{problem}

\textbf{Solucao}:

A concentracao de ${}^{14}$C decai exponencialmente:
\begin{equation}
C(t) = C_0\,e^{-\lambda t}, \qquad \lambda = \frac{\ln 2}{t_{1/2}} = \frac{\ln 2}{5700}\;\text{anos}^{-1}.
\end{equation}

\textbf{1. Concentracao apos $t = 570.000$ anos}:
\begin{align}
n_{1/2} &= \frac{570.000}{5700} = 100 \quad \text{meias-vidas}, \\
C &= C_0\cdot\left(\tfrac{1}{2}\right)^{100} = \frac{1}{10^{12}}\cdot\frac{1}{2^{100}}.
\end{align}
Como $2^{10} \approx 10^3$, temos $2^{100} \approx 10^{30}$, portanto:
\begin{equation}
C \approx \frac{1}{10^{12}} \cdot \frac{1}{10^{30}} = \frac{1}{10^{42}} \quad \text{atomos por atomo de carbono.}
\end{equation}

\textbf{2. Incerteza na idade via propagacao de erros}:

De $t = -\ln(C/C_0)/\lambda$, derivando em relacao a $C$:
\begin{equation}
\delta t = \frac{1}{\lambda}\,\frac{\delta C}{C} = \frac{t_{1/2}}{\ln 2}\,\frac{\delta C}{C}.
\end{equation}
Com $\delta C/C = 10^{-23}$ e $t_{1/2}/\ln 2 = 5700/0{,}693 \approx 8225$ anos:
\begin{equation}
\delta t \approx 8225\;\text{anos} \times 10^{-23} \approx 8 \times 10^{-20}\;\text{anos.}
\end{equation}

\textbf{3. Critica}:

A concentracao esperada ($\sim 10^{-42}$) e 19 ordens de grandeza abaixo da
precisao instrumental citada ($\sim 10^{-23}$). Portanto:
\begin{itemize}
    \item O sinal de ${}^{14}$C seria absolutamente indetectavel apos 570.000 anos.
    \item O metodo ${}^{14}$C so e confiavel ate $\sim 50.000$ anos
    (cerca de 8--9 meias-vidas), pois apos isso a quantidade restante fica abaixo
    do limiar de deteccao.
    \item Qualquer ${}^{14}$C detectado em amostras suposta\-mente mais antigas
    e provavelmente contaminacao moderna ou artefato instrumental.
\end{itemize}

\begin{remark}
Esse problema testa tres competencias do capitulo 5: (1) calcular decaimento
exponencial com meias-vidas, (2) propagar incertezas usando derivadas, (3) interpretar
criticamente os limites do modelo. A conclusao central e que o modelo exponencial
implica um limite pratico de datacao bem antes de 570.000 anos.
\end{remark}

\section{Resumo final}

Comportamento exponencial e um dos blocos mais importantes da modelagem matematica.
Ele aparece quando a dinamica depende do proprio estado e permite previsoes simples e robustas.

Ideias-chave do capitulo 5:
\begin{itemize}
    \item EDO linear de primeira ordem com coeficiente constante;
    \item crescimento e decaimento exponencial;
    \item meia-vida e constante de tempo;
    \item leitura de dados em escala logaritmica;
    \item conexao entre linearizacao local e estabilidade.
\end{itemize}

\end{document}
